\documentclass[11pt]{article}
\ifdefined\pdfmapfile\pdfmapfile{+lm.map}\pdfmapfile{+cm.map}\pdfmapfile{+symbols.map}\fi
\usepackage[T1]{fontenc}
\usepackage{lmodern}
\usepackage[margin=1in]{geometry}
\usepackage{amsmath,amssymb,amsthm,hyperref,listings}
\hypersetup{hidelinks}
\newtheorem{theorem}{Theorem}
\newtheorem{corollary}[theorem]{Corollary}
\newcommand{\LP}{\mathcal{LP}^{+}}
\lstset{basicstyle=\ttfamily\footnotesize,frame=single,breaklines=true,
columns=fullflexible,keepspaces=true,showstringspaces=false}
\title{Counterexamples to Iterated Tur\'an Positivity\\for Multiplier Sequences}
\author{Working proof note}
\date{November 2026}
\begin{document}
\maketitle
\begin{abstract}
For every integer $K\ge6$, the strictly positive multiplier sequence with
exponential generating function $e^x(x+10K)^2(x+K)^2$ violates the sixth
iterated Tur\'an inequality at index $K$. This gives a negative answer to
Problem~3.1 of Csordas and Dimitrov (2000), in its original interior-index
convention. The proof reduces to signs of 131 integer coefficients of a
degree-130 polynomial and includes a complete verifier using only integer
arithmetic. A single example admits a short rational sign certificate.
An explicit finite polynomial with only negative real zeros gives a further
counterexample. No claim of priority or external peer review is made.
\end{abstract}

\section{The original question and the index convention}
For a sequence $\gamma=(\gamma_n)_{n\ge0}$, put
\begin{equation}\label{eq:turan}
T_0(n;\gamma)=\gamma_n,\qquad
T_r(n;\gamma)=T_{r-1}(n;\gamma)^2
 -T_{r-1}(n-1;\gamma)T_{r-1}(n+1;\gamma)
\quad(n\ge r\ge1).
\end{equation}
Csordas and Dimitrov~\cite[Problem~3.1]{CD} ask whether
\begin{equation}\label{eq:question}
T_r(n;\gamma)\ge0\qquad(n\ge r\ge2)
\end{equation}
holds for every nonnegative multiplier sequence. By the P\'olya--Schur
characterization, this is equivalent to requiring the exponential generating
function $\varphi(x)=\sum_{n\ge0}\gamma_nx^n/n!$ to belong to $\LP$.
Here $\LP$ consists of the positive type-I Laguerre--P\'olya functions,
which have a representation
\[
c x^m e^{\sigma x}\prod_j(1+x/a_j),\qquad
c,\sigma\ge0,\quad m\in\mathbb Z_{\ge0},\quad a_j>0,
\quad\sum_j a_j^{-1}<\infty.
\]
Finite products are allowed. In particular, polynomials with nonnegative
coefficients and only nonpositive real zeros belong to this class.

The index restriction $n\ge r$ is essential: it makes every input index
nonnegative without an artificial left-boundary extension. Br\"and\'en's
later restatement~\cite[Problem~8.1]{B} uses the shifted operator
\[
(\mathcal T\gamma)_j=\gamma_{j+1}^2-\gamma_j\gamma_{j+2},\qquad j\ge0.
\]
Induction gives the exact correspondence
\begin{equation}\label{eq:shift}
(\mathcal T^r\gamma)_j=T_r(j+r;\gamma).
\end{equation}
Thus the theorem below answers both formulations. It concerns exponential
generating coefficients; it does not contradict the theorem in~\cite{B}
about ordinary coefficients of polynomials with negative real zeros.
The distinction between these two normalizations is consequential.

The question appears explicitly in the 2000 paper~\cite{CD}, is studied
further in~\cite{CC}, and is restated in~\cite{B}. Wagner's 2022
paper~\cite[\S4.1]{W} still describes infinite log-concavity of multiplier
sequences as conjectural. These sources establish the historical question;
the present mathematical certificate is independent of any priority claim.

\section{A family at every admissible center}
\begin{theorem}\label{thm:family}
For each integer $K\ge6$, let
\[
\varphi_K(x)=e^x(x+10K)^2(x+K)^2
 =\sum_{n=0}^{\infty}\gamma_n(K)\frac{x^n}{n!}.
\]
Then $\gamma(K)$ is a strictly positive multiplier sequence and
\[
\boxed{T_6(K;\gamma(K))<0.}
\]
Moreover, all lower iterates in the dependency window are positive:
\[
T_r(K+c;\gamma(K))>0
\quad\text{if }1\le r\le5\text{ and }|c|\le6-r.
\]
\end{theorem}

\begin{proof}
The factorization
\[
\varphi_K(x)=100K^4e^x(1+x/(10K))^2(1+x/K)^2
\]
is a canonical $\LP$ representation. Its two zeros are $-10K$ and $-K$,
each of multiplicity two. Expanding the polynomial factor gives
\[
(x+10K)^2(x+K)^2=x^4+22Kx^3+141K^2x^2+220K^3x+100K^4.
\]
Writing $(n)_j=n(n-1)\cdots(n-j+1)$, with $(n)_0=1$, series
multiplication therefore gives
\begin{equation}\label{eq:gamma}
\gamma_n(K)=(n)_4+22K(n)_3+141K^2(n)_2+220K^3n+100K^4.
\end{equation}
Every term on the right is nonnegative for an integer $n\ge0$, and the
last term is positive. This also proves strict positivity of the sequence.

To prove the failed inequality, form the following integer polynomials in
the parameter $K$:
\begin{align}
G_{0,c}(K)&=\gamma_{K+c}(K),&&-6\le c\le6,\label{eq:G0}\\
G_{r,c}(K)&=G_{r-1,c}(K)^2
 -G_{r-1,c-1}(K)G_{r-1,c+1}(K),&&|c|\le6-r.
\label{eq:G}
\end{align}
The parameter $K$ is held fixed while the sequence index moves. In
particular, neither adjacent factor in~\eqref{eq:G} has parameter $K-1$
or $K+1$. Induction on $r$ gives
\[
G_{r,c}(K)=T_r(K+c;\gamma(K)).
\]
All its indices are allowed since $K+c\ge K-(6-r)\ge r$.

The complete integer-arithmetic verifier in Appendix~\ref{app:verifier}
constructs the polynomials~\eqref{eq:G0}--\eqref{eq:G} directly. Its output
certifies the following finite identities and coefficient signs:
\begin{enumerate}
\item Every polynomial in row $r$, for $1\le r\le6$, has degree
$2+2^{r+1}$. The row degrees are
\[6,10,18,34,66,130.\]
\item If
\[
G_{6,0}(t+6)=\sum_{i=0}^{130}q_i t^i,
\]
then $q_i<0$ for every $0\le i\le130$.
\item Each of the $11+9+7+5+3=35$ polynomials
$G_{r,c}(t+6)$, where $1\le r\le5$ and $|c|\le6-r$, has every
coefficient strictly positive.
\end{enumerate}
The change of origin in these checks uses only the binomial identity
\[
[t^i]H(t+6)=\sum_{j=i}^{\deg H}[K^j]H(K)\binom ji6^{j-i}.
\]
All input coefficients, polynomial products, and comparisons are integers.
The verifier generates its coefficient tables and does not assume an
externally supplied table is correct. Independent polynomial and direct
integer implementations give the same output.

For $K\ge6$, take $t=K-6\ge0$. The strictly negative constant coefficient
and the nonpositive remaining terms yield $G_{6,0}(K)<0$. The positive
coefficient statements similarly give all asserted lower inequalities.
This proves the theorem for every integer $K\ge6$, rather than just a
finite list of tested parameters.
\end{proof}

The quantifiers should be kept explicit: there is a counterexample sequence
\emph{for each} admissible center $K$. The sequence itself depends on $K$.
The theorem does not say that one fixed sequence fails at every center.

\section{A short certificate for one example}
Taking $K=10$ in~\eqref{eq:gamma} gives
\begin{equation}\label{eq:single}
\boxed{\gamma_n=n^4+214n^3+13451n^2+206334n+1000000.}
\end{equation}
Its generating function is $e^x(x+100)^2(x+10)^2$. This example alone
suffices to answer~\eqref{eq:question} negatively; the following shorter
calculation is independent of the degree-130 certificate.
\begin{corollary}\label{cor:single}
For the sequence~\eqref{eq:single}, $T_5(9),T_5(10),T_5(11)>0$ and
\[
\boxed{-\frac{451}{100000}
<\frac{T_6(10)}{T_5(9)T_5(11)}
<-\frac{450}{100000}.}
\]
Equivalently, $(\mathcal T^6\gamma)_4<0$ in the convention of~\cite{B}.
\end{corollary}
\begin{proof}
Only $\gamma_4,\ldots,\gamma_{16}$ enter the calculation. The row lengths
are $13,11,9,7,5,3,1$, totaling 49 integers. The following exact code
constructs them and proves the displayed rational bounds by integer
cross-multiplication:
\par\noindent
\begin{minipage}{\linewidth}
\begin{lstlisting}
def g(n):
    return n**4+214*n**3+13451*n**2+206334*n+1000000
a = {n:g(n) for n in range(4,17)}
for r in range(1,6):
    a = {n:a[n]**2-a[n-1]*a[n+1]
         for n in range(4+r,17-r)}
    assert min(a.values()) > 0
D = a[9]*a[11]
V = a[10]**2-D
assert D > 0
assert -451*D < 100000*V < -450*D
\end{lstlisting}
\end{minipage}\par
Here $V=T_6(10)$ and $D=T_5(9)T_5(11)$. There is no floating-point
rounding or numerical tolerance in this certificate.
For an alternative compact integer sign certificate, let $S=10^{135}$.
The same recurrence gives
\[
\left\lfloor T_5(9)/S\right\rfloor=10256656,\qquad
\left\lfloor T_5(10)/S\right\rfloor=21535143,\qquad
\left\lfloor T_5(11)/S\right\rfloor=45420371.
\]
Consequently,
\begin{align*}
T_5(10)^2/S^2 &< 21535144^2=463762427100736\\
&<465861120739376=10256656\cdot45420371\\
&\le T_5(9)T_5(11)/S^2,
\end{align*}
which directly proves $T_6(10)<0$.

\end{proof}

\section{Finite polynomials also give counterexamples}
\begin{corollary}\label{cor:finite}
For each fixed integer $K\ge6$, all sufficiently large integers $N$ give a
counterexample with polynomial exponential generating function
\[
P_{K,N}(x)=(1+x/N)^N(x+10K)^2(x+K)^2.
\]
In particular, $K=10$ and $N=50000$ work. Write
\[\eta_n=n![x^n]P_{10,50000}(x).\]
Then
\[
\boxed{-\frac{2}{1000}
<\frac{T_6(10;\eta)}{T_5(9;\eta)T_5(11;\eta)}
<-\frac{1}{1000}.}
\]
The denominator is positive.
\end{corollary}
\begin{proof}
All zeros of $P_{K,N}$ are negative real numbers. For fixed $K$, its
coefficients at any fixed finite set of indices converge to those of
$\varphi_K$ as $N\to\infty$. For example,
$[x^m](1+x/N)^N=\binom Nm N^{-m}\to1/m!$.
The expression $T_6(K)$ is a polynomial in the thirteen coefficients with
indices $K-6,\ldots,K+6$, so it is continuous in those coefficients.
The strict negative limit in Theorem~\ref{thm:family} proves the claim for
all sufficiently large $N$; this argument alone gives no explicit threshold.

For the concrete value $N=50000$, use the positive geometric rescaling
$a_n=N^n\eta_n$. Induction in~\eqref{eq:turan} gives
\[
T_r(n;a)=N^{2^r n}T_r(n;\eta).
\]
Both terms in the recurrence have the same scaling exponent, so signs and
the displayed normalized ratio are unchanged. Put
$(p_0,p_1,p_2,p_3,p_4)=(1000000,220000,14100,220,1)$. Then
\[
a_n=\sum_{j=0}^{\min(4,n)}p_j(n)_j(N)_{n-j}N^j
\]
is an integer. Appendix~\ref{app:finite} evaluates these integers and
verifies the rational enclosure. The polynomial has degree $50004$.
\end{proof}

\section{Scope and reproducibility}
The result is a counterexample to the universal assertion in the original
2000 problem, within its original domain $n\ge r$. No padding convention
at negative indices is involved. The first five iterates are certified
positive on the dependency windows stated in Theorem~\ref{thm:family};
we do not assert their positivity at every index, nor that iteration six
is the smallest possible order of failure.

The result has no implication that the Riemann hypothesis is false:
the generating functions here are the elementary functions displayed above.
It also does not disprove eventual-positivity conjectures for each fixed
sequence, whose quantifiers differ from those of Theorem~\ref{thm:family}.

The full verifier creates \texttt{turan\_theorem\_certificate.json},
containing all 131 polynomial coefficients before and after the shift,
and all 49 integers of the single-example recurrence. A separately written
symbolic-polynomial implementation reproduced the family certificate;
separately written recursive and triangular integer implementations
reproduced the single example. The proofs use exact arithmetic throughout.

A bounded literature search through November 2026 found no earlier
resolution of this exact interior-index question. This is not an exhaustive
priority determination. The proof note is a reproducible mathematical
certificate, not a claim of independent external acceptance.

\appendix
\section{Complete integer verifier for the family}\label{app:verifier}
Run the following file with standard Python 3. No external libraries,
downloaded coefficient tables, floating-point arithmetic, or search
procedure are required. The program reconstructs its entire certificate.
\begin{lstlisting}
#!/usr/bin/env python3
"""Exact certificate: an iterated Turan counterexample at every center K>=6.

Uses only Python's standard library. Coefficient lists are in increasing
powers. The parameter K stays fixed when the sequence index changes.
The coefficient tables are generated, not read as trusted inputs.
"""
from math import comb, prod
from pathlib import Path
import json


def add(a, b):
    out = [0] * max(len(a), len(b))
    for i, c in enumerate(a):
        out[i] += c
    for i, c in enumerate(b):
        out[i] += c
    while len(out) > 1 and out[-1] == 0:
        out.pop()
    return out


def mul(a, b):
    out = [0] * (len(a) + len(b) - 1)
    for i, c in enumerate(a):
        for j, d in enumerate(b):
            out[i+j] += c*d
    while len(out) > 1 and out[-1] == 0:
        out.pop()
    return out


def sub(a, b):
    return add(a, [-c for c in b])


def evaluate(a, x):
    out = 0
    for c in reversed(a):
        out = out*x + c
    return out


def shift(a, h):
    return [sum(a[j]*comb(j, i)*h**(j-i)
                for j in range(i, len(a))) for i in range(len(a))]


def input_polynomial(c):
    # Coefficients of (x+10K)^2(x+K)^2 in powers of x.
    p = ([0,0,0,0,100], [0,0,0,220], [0,0,141], [0,22], [1])
    falling, ans = [1], [0]
    for j in range(5):
        ans = add(ans, mul(p[j], falling))
        falling = mul(falling, [c-j, 1])
    return ans


def gamma(j, K):
    p = (100*K**4, 220*K**3, 141*K**2, 22*K, 1)
    return sum(p[r]*prod(range(j-r+1, j+1))
               for r in range(min(j, 4)+1))


def integer_rows(K):
    row = [gamma(j, K) for j in range(K-6, K+7)]
    rows = [row]
    for r in range(6):
        row = [row[i]**2-row[i-1]*row[i+1]
               for i in range(1, len(row)-1)]
        rows.append(row)
    return rows


def verify():
    row = [input_polynomial(c) for c in range(-6, 7)]
    degrees = []
    lower_count = 0
    for r in range(1, 7):
        row = [sub(mul(row[i], row[i]), mul(row[i-1], row[i+1]))
               for i in range(1, len(row)-1)]
        degree_set = sorted(set(len(p)-1 for p in row))
        assert degree_set == [2+2**(r+1)]
        degrees.append(degree_set[0])
        if r <= 5:
            for p in row:
                assert all(c > 0 for c in shift(p, 6))
                lower_count += 1
    assert len(row) == 1 and len(row[0]) == 131
    coeff = row[0]
    shifted = shift(coeff, 6)
    assert all(c < 0 for c in shifted)
    assert lower_count == 35

    # Direct checks hold K fixed and move only the sequence index.
    for K in (6, 10, 37):
        rows = integer_rows(K)
        assert rows[-1][0] == evaluate(coeff, K)

    # A compact certificate for the single example K=10.
    rows = integer_rows(10)
    for j, value in zip(range(4, 17), rows[0]):
        assert value == j**4+214*j**3+13451*j**2+206334*j+1000000
    assert all(c > 0 for row in rows[:-1] for c in row)
    left, middle, right = rows[5]
    D = left*right
    V = middle**2-D
    assert V == rows[6][0] and D > 0
    assert -451*D < 100000*V < -450*D
    return {
        'family': 'exp(x)*(x+10*K)^2*(x+K)^2; integer K>=6',
        'degrees': degrees,
        'negative_coefficients_after_K_equals_t_plus_6': len(shifted),
        'positive_lower_window_polynomials': lower_count,
        'coefficients_in_K': coeff,
        'coefficients_in_K_minus_6': shifted,
        'K10_rows_starting_at_indices_4_through_10': rows,
        'K10_T6': V,
        'K10_ratio_denominator': D,
        'K10_ratio_bounds': ['-451/100000', '-450/100000'],
    }


if __name__ == '__main__':
    report = verify()
    path = Path(__file__).with_name('turan_theorem_certificate.json')
    path.write_text(json.dumps(report, indent=2)+'\n')
    print('PASS: all 131 shifted coefficients are strictly negative.')
    print('PASS: all 35 lower-window polynomials have positive coefficients.')
    print('PASS: direct fixed-parameter checks at K=6, 10, 37 agree.')
    print('PASS: exact K=10 ratio lies between -0.00451 and -0.00450.')
    print('Generated certificate:', path)
\end{lstlisting}

\clearpage
\section{Exact verifier for the finite polynomial}\label{app:finite}
The following independent calculation uses only integer arithmetic and
exact rational comparisons.
\begin{lstlisting}
#!/usr/bin/env python3
"""Finite-polynomial multiplier-sequence counterexample, exact integers only.

Phi(x)=(1+x/50000)^50000*(x+100)^2*(x+10)^2.
All zeros are strictly negative. With gamma_j=j![x^j]Phi, we verify
T_6(10)<0 after the harmless geometric rescaling a_j=50000^j*gamma_j.
"""
from math import prod
from fractions import Fraction

N=50000
P=(1000000,220000,14100,220,1)
def falling(n,r):
    return prod(range(n-r+1,n+1)) if r<=n else 0

def coefficient(j):
    return sum(c*falling(j,r)*falling(N,j-r)*N**r
               for r,c in enumerate(P) if r<=j)

def verify():
    row=[coefficient(j) for j in range(4,17)]
    for r in range(1,6):
        row=[row[i]**2-row[i-1]*row[i+1] for i in range(1,len(row)-1)]
        if not all(c>0 for c in row):
            raise ArithmeticError('A lower dependency-window term is not positive')
    denominator=row[0]*row[2]
    last=row[1]**2-denominator
    if last>=0:raise ArithmeticError('Finite-polynomial counterexample failed')
    ratio=Fraction(last,denominator)
    # Small rational enclosure, independently meaningful without long integers.
    if not Fraction(-2,1000)<ratio<Fraction(-1,1000):
        raise ArithmeticError('Stated rational enclosure failed')
    print('PASS: degree50004 polynomial with only negative zeros has T6(10)<0.')
    print('Exact ratio T6(10)/(T5(9)*T5(11)) lies strictly between -2/1000 and -1/1000.')

if __name__=='__main__':verify()
\end{lstlisting}

{\small
\begin{thebibliography}{9}
\bibitem{CD}
G.~Csordas and D.~K.~Dimitrov,
\emph{Conjectures and theorems in the theory of entire functions},
Numerical Algorithms \textbf{25} (2000), 109--122.
\href{https://doi.org/10.1023/A:1016604906346}{doi:10.1023/A:1016604906346}.
Author manuscript, Problem~3.1:
\url{https://www.dcce.ibilce.unesp.br/~dimitrov/papers/vargana.pdf}.
\bibitem{CC}
T.~Craven and G.~Csordas,
\emph{Iterated Laguerre and Tur\'an inequalities},
Journal of Inequalities in Pure and Applied Mathematics \textbf{3}(3)
(2002), Article~39.
\url{https://math.hawaii.edu/~tom/mathfiles/iteratedLT.pdf}.
\bibitem{B}
P.~Br\"and\'en,
\emph{Iterated sequences and the geometry of zeros},
Journal f\"ur die reine und angewandte Mathematik \textbf{658}
(2011), 115--131.
\href{https://doi.org/10.1515/CRELLE.2011.063}{doi:10.1515/CRELLE.2011.063}.
\url{https://arxiv.org/abs/0909.1927}; arXiv version 3, \S8, Problem~8.1.
\bibitem{W}
I.~Wagner,
\emph{On a new class of Laguerre--P\'olya type functions with applications
in number theory}, Pacific Journal of Mathematics \textbf{320}(1)
(2022), 177--192.
\href{https://doi.org/10.2140/pjm.2022.320.177}{doi:10.2140/pjm.2022.320.177}.
\url{https://arxiv.org/abs/2108.01827}; arXiv version 2, \S4.1, p.~10.
\end{thebibliography}
}
\end{document}
